\section{Limitations and reproducibility}
The dataset is retrospective and reflects its source population and measurement practices. Small five-fold samples, especially for the less prevalent outcomes, produce variable estimates; development CV summaries and bootstrap intervals are descriptive, not guarantees of deployment performance. The threshold constraint of 0.80 sensitivity is an operating-point assumption. T06 notes that the full-development estimator is trained on more records than the outer-fold estimators used to produce calibration scores, so calibration transfer can be imperfect. The FIBR\_PREDS Platt mapping slightly worsened development Brier and log loss relative to raw scores, and the held-out FIBR ranking/calibration results are weaker than for the other outcomes.

The most important alignment limitation is preserved from T07: the Stage~1 source uses a shared \texttt{test\_idx} to create X, y, and IDs, and the saved test row count is checked. However, the existing CSVs contain no shared row-level key or alignment fingerprint. Thus exact row-by-row correspondence between saved test predictors and IDs cannot be independently proven from current artifacts. The split was not recreated. This limitation does not change the recorded evaluation, but it should accompany interpretation of row-specific audit outputs.

The T07 report includes evaluation labels for audit; the T08 submission contains only IDs and probabilities. The mortality SHAP explanation uses development data and is associational. This draft leaves all six assignment answers unfinished. The optional local SHAP explanations for two patients are also not included.

Reproducibility sources include the versioned scripts and run manifests: Stage~1 split and feature creation; Stage~6 common-fold model comparison; T06 frozen model manifest, bundles, thresholds, OOF predictions, and freeze report; T07 report/manifest/aggregate tables; T08 provenance manifest and CSV; and T09 SHAP script/manifest. The project environment and commands are documented in the run logs. No training or evaluation was rerun to produce this report. See \texttt{README.md} for the artifact map.
